\begin{table}[htbp]
\centering
\caption{P-Tree model performance when the number of boosting iterations is set to 2 (\texttt{num\_iter = 2}). The in-sample results are better but the out-of-sample results are worse. This shows that deeper trees tend to overfit as they capture noise rather than true relationships.}
\label{tab:performance_2splits}
\setlength{\tabcolsep}{3pt}
\renewcommand{\arraystretch}{1.0}
\small
\begin{tabular}{lccccccc}
\hline
Scenario & Ret (\%) & Vol (\%) & Sharpe & $\beta$ & CAPM $\alpha$ (\%, t) & FF3 $\alpha$ (\%, t) & $R^2$ (\%) \\
\hline
Market (VW) & 10.9 & 19.0 & 0.64 & 1.00 & -- & -- & -- \\
A & 10.9 & 9.7 & 1.12 & 0.07 & 8.19*** (4.21) & 7.84*** (4.26) & -- \\
B (Train) & 13.2 & 11.4 & 1.15 & -0.05 & 13.55*** (5.00) & 13.55*** (4.95) & -- \\
B (Test) & 1.3 & 5.8 & 0.22 & 0.02 & 0.35 (0.27) & 0.45 (0.37) & 0.3 \\
C (Train) & 9.4 & 5.2 & 1.83 & 0.04 & 9.10*** (5.10) & 9.17*** (5.38) & -- \\
C (Test) & 3.4 & 6.8 & 0.50 & -0.02 & 3.11 (1.00) & 2.71 (0.85) & 0.5 \\
\hline
\end{tabular}

\vspace{0.15cm}
\begin{minipage}{0.9\textwidth}
\footnotesize
\textit{Notes:} Ret and Vol are annualized. Alphas are annualized (\%) with Newey--West t-stats (12 lags). Significance: *** $p<0.01$, ** $p<0.05$, * $p<0.10$. $R^2$ is out-of-sample.
\end{minipage}
\end{table}
